\subsection{Level-Set Normalization and the Improved Empirical Fisher}

For the loss-functional specialization $\phi_a=\ell_A$ and Euclidean output metric $G_i=I$, the matrix from the previous section is the improved empirical Fisher (iEF) introduced by Wu et al. \citep{IEF}:
\[
M_{2,I}=F_{\mathrm{IEF}}.
\]
Hence iEF is the quadratic form of a first-order local constraint that protects old behavior geometrically: instead of preserving a raw scalar-loss coordinate, it penalizes only output motion normal to the old loss level set. The original iEF paper derives the same normalization for approximate natural-gradient descent. In that setting, the observed-label OPG reweights sample contributions by squared output-gradient magnitude; iEF removes this scale factor and retains the normalized output-loss direction. Here the normalization arises from metric projection onto the normal quotient of the old-loss tangent space.

Consider the scalar least-squares setting in Euclidean geometry. Suppose the model output is scalar, such that
\[
\ell_n(z)=\frac12(z-y_n)^2.
\]
Then
\[
r_n=z_n^A-y_n
\]
is scalar, and for $r_n\neq 0$ we have
\[
\frac{r_n^2}{r_n^2}=1.
\]
Therefore the \emph{undamped} Euclidean IEF contribution reduces to
\[
J_n^\top J_n,
\]
and averaging gives
\[
M_{\mathrm{IEF}}
=
\frac{1}{N}\sum_{n=1}^N J_n^\top J_n.
\]
This is exactly the Gauss--Newton or output-Jacobian form for squared-error regression \citep{GaussNewton,MartensHF,MartensNG}, the standard practical preconditioning matrix for second-order optimization. This matches our geometric intuition: loss preservation is governed by how parameters move the output, which is encoded by the Jacobian $J$, not by the current output-gradient magnitude. This equivalence was made explicit in \citep{IEF}.
